\begin{definition}
    \textbf{Equipotenza}

    Due insiemi $\mathbb{X}$ e $\mathbb{Y}$ si dicono equipotenti se esiste una funzione
    biunivoca dal primo al secondo o, equivalentemente (in quanto biunivoca),
    dal secondo al primo. Affermazione equivalente ad affermare l'uguaglianza delle
    cardinalità dei due insiemi.
    
    In simboli:
    \[
        \left|\mathbb{X}\right| = \left|\mathbb{Y}\right| \Leftrightarrow
        \exists f:\mathbb{X}\rightarrow\mathbb{Y}\,\big|
        \begin{cases}
            \forall x_1, x_2 \in\mathbb{X}, x_1\neq x_2 \rightarrow f\left(x_1\right)\neq f\left(x_2\right) \quad \text{($f$ inj.)}\\
            \forall y \in\mathbb{Y}, \exists x\in\mathbb{X} \mid y=f\left(x\right) \quad \text{($f$ sur.)}
        \end{cases}.
    \]
\end{definition}

Si può dunque facilmente dimostrare, grazie alla condizione
$\left|\mathbb{X}\right| = \left|\mathbb{Y}\right|$, che la relazione di equipotenza, indicata
d'ora in avanti con $\eqpt$ sia una relazione d'equivalenza.

\begin{proposition}
    La relazione di equipotenza $\eqpt$ è una relazione d'equivalenza.
\end{proposition}

\begin{definition}
    \textbf{Insieme finito}

    Sia $\mathbb{X}$ un qualsiasi insieme, si dice che $\mathbb{X}$ è un insieme finito
    se e solo se
    \[
        \exists n \in\mathbb{N} \mid X\eqpt\left\{1,2,3,\ldots, n\right\}.
    \]
\end{definition}

\begin{definition}
    \textbf{Insieme infinito}

    Sia $\mathbb{X}$ un qualsiasi insieme, si dice che $\mathbb{X}$ è un insieme infinito
    se e solo se non è finito.
\end{definition}

\begin{definition}
    \textbf{Insieme numerabile}

    Dato un insieme $\mathbb{X}$ infinito, si dice che $\mathbb{X}$ è numerabile
    se e solo se $\mathbb{X}$ è equipotente a $\mathbb{N}$, o, equivalentemente,
    se $\mathbb{N}$ è equipotente a $\mathbb{X}$.
    
    In simboli:

    \[ \mathbb{N} \eqpt \mathbb{X}. \]

    \leavevmode

    Deve dunque esistere una funzione biunivoca definita come $f:\mathbb{N}\rightarrow\mathbb{X}$;
    questa funzione $f$, avendo come dominio $\mathbb{N}$, per \defref{def:successione}
    risulta quindi una serie. Di conseguenza possiamo dedurre che se un qualsiasi insieme
    infinito $\mathbb{X}$ è numerabile allora possiamo rappresentarlo come una serie
    $\mathbb{X}=\left\{x_n\right\}_{n\in\mathbb{N}}$.
\end{definition}

\begin{proposition}\label{pro:neqpt-fin}
    Siano $\mathbb{X}_1$ e $\mathbb{X}_2$ due insiemi tali che $\mathbb{X}_1 \subsetneq \mathbb{X}_2$,
    possiamo affermare che $\mathbb{X}_1 \neqpt \mathbb{X}_2$ se e solo se almeno uno dei due insiemi
    è finito.
\end{proposition}

\begin{example}
    (\textit{di} \proref{pro:neqpt-fin})

    Prendiamo la sequente legge per $n\in\mathbb{N}_0$
    \[ n\mapsto n+1.\]
    La legge sopra definita implica l'esistenza di una funzione $f$ definita come segue
    \[ f:\mathbb{N}_0 \rightarrow \mathbb{N}, \qquad n\mapsto n+1, \]
    questa funzione $f$ è biunivoca, dunque implica che
    \[ \mathbb{N} \eqpt \mathbb{N}_0, \]
    nonostante
    \[ \mathbb{N} \subsetneq \mathbb{N}_0. \]

    Un insieme infinito può quindi essere equipotente a un suo sottoinsieme infinito.
\end{example}

\begin{example}
    (\textit{di} \proref{pro:neqpt-fin})

    Prendiamo la sequente legge per $n\in\mathbb{N}$
    \[ n\mapsto 2n\]
    e introduciamo la seguente notazione
    \[2\mathbb{N}:=\left\{2n\mid n\in\mathbb{N}\right\}.\]
    La legge sopra definita implica l'esistenza di una funzione $f$ definita come segue
    \[ f:\mathbb{N} \rightarrow 2\mathbb{N}, \qquad n\mapsto 2n, \]
    questa funzione $f$ è biunivoca, dunque implica che
    \[ 2\mathbb{N} \eqpt \mathbb{N}, \]
    nonostante
    \[ 2\mathbb{N} \subsetneq \mathbb{N}. \]

    Un insieme infinito può quindi essere equipotente a un suo sottoinsieme infinito.
\end{example}

\begin{theorem}
    \textbf{Unione tra numerabile e finito}

    Dato un insieme $\mathbb{X}$ equipotente a $\mathbb{N}$ e un insieme $\mathbb{A}$ finito,
    allora anche l'unione tra $\mathbb{X}$ e $\mathbb{A}$ sarà equipotente a $\mathbb{N}$

    In simboli:
    
    \[
        \exists n \in\mathbb{N} \mid
            \mathbb{X} \eqpt \mathbb{N} \land
            \mathbb{A}\eqpt\left\{0,1,2,\ldots, n\right\}
        \implies \mathbb{X} \cup\mathbb{A}\eqpt\mathbb{N}.
    \]
\end{theorem}

\begin{proposition}
    Il numerabile è la "più piccola" cardinalità degli infiniti.
\end{proposition}

\begin{proposition}
    Per ogni insieme infinito esiste un suo sottoinsieme numerabile.
    Sia $\mathbb{X}$ un qualsiasi insieme infinto, allora in simboli abbiamo che
    \[\forall \mathbb{X}, \exists \mathbb{A}\subset\mathbb{X} \mid \mathbb{A}\eqpt\mathbb{N}.\]
\end{proposition}

\begin{proposition}\label{pro:sot-inf-num}
    Un sottoinsieme infinito di un insieme numerabile è a sua volta numerabile.

    In simboli:
    
    
    \[
        \nexists n \in\mathbb{N} \mid
            \mathbb{X} \eqpt \mathbb{N} \land
            \mathbb{A} \subseteq \mathbb{X} \land
            \mathbb{A}\eqpt\left\{0,1,2,\ldots, n\right\}
        \implies \mathbb{A}\eqpt\mathbb{N}.
    \]
\end{proposition}

\begin{demonstration}
    \textit{di} \proref{pro:sot-inf-num}

    \[ \Hp \mathbb{N}\eqpt\mathbb{X} \;\;\land\;\; \mathbb{A} \subset \mathbb{X} \;\;\land\;\; \forall n\in\mathbb{N}, \mathbb{A}\neqpt \left\{m\in\mathbb{N}\mid m\leq n\right\}\]
    \[ \Th \mathbb{A}\eqpt\mathbb{N} \]

    \begin{enumerate}[label=(\roman*)]
        \item
            \[ \mathbb{N}\eqpt\mathbb{X}, \]
            \[ \exists f:\mathbb{N}\rightarrow\mathbb{X} \mid f \text{ è biettiva}, \]
            \[ \exists f^{-1}:\mathbb{X}\rightarrow\mathbb{N} \mid f^{-1} \text{ è biettiva}. \]

        \item Definiamo la seguente funzione
        \[ g:\mathbb{N} \rightarrow \mathbb{A}, \]
        \[ g\left(1\right) := f\left(\min \left\{p\in\mathbb{N}\mid f\left(p\right)\in\mathbb{A}\right\}\right), \]

        \[ \forall m > 1, \;\;g\left(m\right) := f\left( \min \left\{p\in\mathbb{N}\mid f\left(p\right)\in\mathbb{X} \land p > f^{-1}\left(g\left(m-1\right)\right)\right\}\right). \]

        \item Mostriamo che $g$ è iniettiva
        \begin{enumerate}[label=\arabic*.]
            \item Ragioniamo $\forall m \in\mathbb{N}\setminus\left\{1\right\}$.
            \item Essendo sia $f$ che $f^{-1}$ biettive, allora
                \[ g\left(m\right) = f\left( \min \left\{p\in\mathbb{N}\mid f\left(p\right)\in\mathbb{X} \land p > f^{-1}\left(g\left(m-1\right)\right)\right\}\right), \]
                \[ f^{-1}\left(g\left(m\right)\right) = f^{-1}\left(f\left( \min \left\{p\in\mathbb{N}\mid f\left(p\right)\in\mathbb{X} \land p > f^{-1}\left(g\left(m-1\right)\right)\right\}\right)\right), \]
                \[ f^{-1}\left(g\left(m\right)\right) = \min \left\{p\in\mathbb{N}\mid f\left(p\right)\in\mathbb{X} \land p > f^{-1}\left(g\left(m-1\right)\right)\right\}.\]

            \item Prendendo $h:f^{-1}\circ g$,
                \[ h\left(m\right) = \min \left\{p\in\mathbb{N}\mid f\left(p\right)\in\mathbb{X} \land p > h\left(m-1\right)\right\},\]
                \[ h\left(m\right) \in \left\{p\in\mathbb{N}\mid f\left(p\right)\in\mathbb{X} \land p > h\left(m-1\right)\right\},\]
                \[ h\left(m\right) \in\mathbb{N} \land f\left(h\left(m\right)\right)\in\mathbb{X} \land h\left(m\right) > h\left(m-1\right) .\]

            \item Abbiamo dunque che
                \[ h\left(m\right) > h\left(m-1\right), \]
                \[ h\left(m\right) > h\left(m-1\right) > h\left(m-2\right) > \ldots > h\left(2\right). \]
            \item Sapendo che per $m=2$, allora $h\left(2\right)>h\left(1\right)$\label{pas:ord}
                \[ h\left(m\right) > h\left(m-1\right) > h\left(m-2\right) > \ldots > h\left(1\right), \]
                \[ \forall m_1, m_2\in\mathbb{N}, \;\;m_1>m_2 \rightarrow h\left(m_1\right) > h\left(m_2\right). \]

            \item Dunque, ragionando $\forall m_1, m_2\in\mathbb{N}$,
                \[
                    m_1 \neq m_2,\]\[
                    m_1 > m_2 \lor m_1<m_2 ,\]\[
                    h\left(m_1\right) > h\left(m_2\right)\lor h\left(m_1\right) < h\left(m_2\right) ,\]\[
                    h\left(m_1\right) \neq h\left(m_2\right)\lor h\left(m_1\right) \neq h\left(m_2\right) ,\]\[
                    h\left(m_1\right) \neq h\left(m_2\right), \]\[
                    f^{-1}\left(g\left(m_1\right)\right) \neq f^{-1}\left(g\left(m_2\right)\right), \]\[
                    f\left(f^{-1}\left(g\left(m_1\right)\right)\right) \neq f\left(f^{-1}\left(g\left(m_1\right)\right)\right), \]\[
                    g\left(m_1\right) \neq g\left(m_1\right).
                \]

            \item In sintesi  
                \[
                    \forall m_1, m_2\in\mathbb{N}, \;\; m_1 \neq m_2 \implies g\left(m_1\right) \neq g\left(m_1\right).
                \]

            

        \end{enumerate}
        \qed

        \item Mostriamo che $g$ è suriettiva
        \begin{enumerate}[label=\arabic*.]
            \item Affinche la funzione sia suriettiva deve valere che
                \[ \forall a\in\mathbb{A}, \exists m\in\mathbb{N}\mid g\left(m\right)=a, \]
                \[ \forall a\in\mathbb{A}, \exists m\in\mathbb{N}\mid f^{1}\left(g\left(m\right)\right)=f^{-1}\left(a\right). \]
                \item Con $h: f^{-1} \circ g$,
                \[ \forall a\in\mathbb{A}, \exists m\in\mathbb{N}\mid h\left(m\right) = f^{-1}\left(a\right). \]
            \item Supponiamo per assurdo che
                \[ \forall a\in\mathbb{A}, \forall m\in\mathbb{N}, h\left(m\right) \neq f^{-1}\left(a\right), \]
                \[ a\in\mathbb{A} \land m\in\mathbb{N} \Rightarrow h\left(m\right) \neq f^{-1}\left(a\right), \]
                \[ a\in\mathbb{A} \land m\in\mathbb{N} \Rightarrow h\left(m\right) \neq f^{-1}\left(a\right). \]
                \[ a\in\mathbb{A} \Rightarrow h\left(1\right) \neq f^{-1}\left(a\right) \land h\left(2\right) \neq f^{-1}\left(a\right) \land \dotsm \land h\left(m\right) \neq f^{-1}\left(a\right) \land \dotsm . \]
            \item Ma sapendo (per \pasref{pas:ord}) che $m_1 > m_2 \implies h\left(m_1\right) > h\left(m_1\right)$
                \[ a\in\mathbb{A} \Rightarrow f^{-1}\left(a\right) < h\left(1\right) \lor h\left(1\right) < f^{-1}\left(a\right) < h\left(2\right) \lor \]
                \[\lor \dotsm \lor h\left(m\right) < f^{-1}\left(a\right) < h\left(m+1\right) \lor \dotsm, \]
                \[ a\in\mathbb{A} \land m\in\mathbb{N} \Rightarrow f^{-1}\left(a\right) < h\left(1\right) \lor h\left(m\right) < f^{-1}\left(a\right) < h\left(m+1\right), \]
                \[ a\in\mathbb{A} \land m\in\mathbb{N} \Rightarrow 0<f^{-1}\left(a\right) < h\left(1\right) \lor h\left(m\right) < f^{-1}\left(a\right) < h\left(m+1\right), \]

            \item Mostriamo che $0<f^{-1}\left(a\right) < h\left(1\right)$ non può essere vero
            \begin{enumerate}[label=\arabic{enumii}.\arabic*.]
                \item Supponiamo per assurdo che $0<f^{-1}\left(a\right) < h\left(1\right)$.
                \item Allora
                    \[ h\left(1\right)=f^{-1}\left(g\left(1\right)\right)= \]
                    \[ =f^{-1}\left(f\left(\min \left\{p\in\mathbb{N}\mid f\left(p\right)\in\mathbb{A}\right\}\right)\right)= \]
                    \[ =\min \left\{p\in\mathbb{N}\mid f\left(p\right)\in\mathbb{A}\right\}. \]
                \item Dunque
                    \[ h\left(1\right) = \min \left\{p\in\mathbb{N}\mid f\left(p\right)\in\mathbb{A}\right\}, \]
                    \[ \forall t \in \left\{p\in\mathbb{N}\mid f\left(p\right)\in\mathbb{A}\right\}, h\left(1\right)\leq t.\]
                \item Ma
                    \[
                        \begin{cases}
                            f^{-1}\left(a\right)\in\mathbb{N} \\
                            f\left(f^{-1}\left(a\right)\right)=a\in\mathbb{A}
                        \end{cases} \Rightarrow f^{-1}\left(a\right)\in\left\{p\in\mathbb{N}\mid f\left(p\right)\in\mathbb{A}\right\}\Rightarrow
                    \]\[\Rightarrow h\left(1\right) \leq f^{-1}\left(a\right)\Rightarrow f^{-1}\left(a\right) < h \left(1\right) \leq f^{-1}\left(a\right) \Rightarrow
                    \]\[
                        \Rightarrow f^{-1}\left(a\right) < f^{-1}\left(a\right).
                    \]

                \item La negaizone della tesi porta a una contraddizione.
                \item La tesi secondo cui $0<f^{-1}\left(a\right) < h\left(1\right)$ non può essere vero è necessariamente vera.
                
            \end{enumerate}
            \qed

            \item Essendo $0<f^{-1}\left(a\right) < h\left(1\right)$ sempre falso, allora $\forall a\in\mathbb{A}, m\in\mathbb{N}$
                \[ 0<f^{-1}\left(a\right) < h\left(1\right) \lor h\left(m\right) < f^{-1}\left(a\right) < h\left(m+1\right), \]
                \[ h\left(m\right) < f^{-1}\left(a\right) < h\left(m+1\right), \]
                \[
                    \begin{cases}
                        h\left(m\right) < f^{-1}\left(a\right) \\
                        f^{-1}\left(a\right) < h\left(m+1\right)
                    \end{cases},
                \]
                \[
                    \begin{cases}
                        h\left(m\right) < f^{-1}\left(a\right) \\
                        f^{-1}\left(a\right) < f^{-1}\left(f\left( \min \left\{p\in\mathbb{N}\mid f\left(p\right)\in\mathbb{X} \land p > f^{-1}\left(g\left(m\right)\right)\right\}\right)\right)
                    \end{cases},
                \]
                \[
                    \begin{cases}
                        h\left(m\right) < f^{-1}\left(a\right) \\
                        f^{-1}\left(a\right) < \min \left\{p\in\mathbb{N}\mid f\left(p\right)\in\mathbb{X} \land p > h\left(m\right)\right\}
                    \end{cases},
                \]
            \item Prendendo $n:=f^{-1}\left(a\right)$
                \[
                    \begin{cases}
                        h\left(m\right) < n \\
                        n < \min \left\{p\in\mathbb{N}\mid f\left(p\right)\in\mathbb{X} \land p > h\left(m\right)\right\}
                    \end{cases},
                \]
                \[
                    \begin{cases}
                        n\in\mathbb{N} \land f\left(n\right)\in\mathbb{X} \land n > h\left(m\right) \\
                        n < \min \left\{p\in\mathbb{N}\mid f\left(p\right)\in\mathbb{X} \land p > h\left(m\right)\right\}
                    \end{cases},
                \]
                \[
                    \begin{cases}
                        n\in\left\{p\in\mathbb{N}\mid f\left(p\right)\in\mathbb{X} \land p > h\left(m\right)\right\} \\
                        n < \min \left\{p\in\mathbb{N}\mid f\left(p\right)\in\mathbb{X} \land p > h\left(m\right)\right\}
                    \end{cases}.
                \]

            \item Prendendo $k\left(m\right) := \min \left\{p\in\mathbb{N}\mid f\left(p\right)\in\mathbb{X} \land p > h\left(m\right)\right\}$, allora
                \[
                    \begin{cases}
                        n\in\left\{p\in\mathbb{N}\mid f\left(p\right)\in\mathbb{X} \land p > h\left(m\right)\right\} \\
                        n < k\left(m\right)\\
                        \forall b\in \left\{p\in\mathbb{N}\mid f\left(p\right)\in\mathbb{X} \land p > h\left(m\right)\right\}, k\left(m\right) \leq b
                    \end{cases},
                \]
                \[
                    \begin{cases}
                        k\left(m\right) \leq n\\
                        n < k\left(m\right)\\
                    \end{cases},
                \]
                \[ n < k\left(m\right) \leq n, \]
                \[ n < n ,\]
                \[ n \neq n .\]
            \item La negazione della tesi porta a una contraddizione.
            \item La tesi secondo cui $g$ è suriettiva è necessariamente vera.
        \end{enumerate}
        \qed

        \item La funzione così $g$ definita è quindi sia iniettiva che suriettiva.
        \item La funzione $g$ è biettiva.
        \item \[ \exists g:\mathbb{N}\rightarrow\mathbb{A} \mid g \text{ è biettiva}, \]
            \[ \mathbb{N}\eqpt\mathbb{A}. \]
    \end{enumerate}
    \qed
\end{demonstration}